
% ============================================================
% SECTION 5: COMBINED BACKEND ARCHITECTURE
% ============================================================
\section{Combined Backend Architecture}

\subsection{Purpose and Design}

The Combined Backend is a design pattern unique to the MASAREN platform that merges all 35+ microservices into a single FastAPI process for local development efficiency. This approach eliminates the overhead of running 30+ separate Docker containers during development while maintaining complete service compatibility.

\begin{description}
  \item[Purpose:] Single-process backend that merges all microservices for local development
  \item[Architecture:] Imports each service as a uniquely-named package and routes requests to appropriate handlers
  \item[Benefits:] Local development on a laptop or small cloud droplet without running the full Docker stack
  \item[Complexity:] Duplicate table registration requires SQLAlchemy Table.\_\_new\_\_ patching
  \item[Performance:] Single process eliminates network overhead between services
\end{description}

\subsection{Service Registration Mechanism}

The combined backend programmatically registers each service package at runtime. Each service is given a unique alias and its source directory is added to Python's module path:

\begin{lstlisting}[language=Python, caption=Service Registration in combined/main.py]
SERVICES = [
    ("ms_identity",      "identity/identity-service"),
    ("ms_auth",          "identity/authentication-service"),
    ("ms_authz",         "identity/authorization-service"),
    ("ms_organization",  "execution/organization-service"),
    ("ms_program",       "execution/program-service"),
    ("ms_task",          "execution/task-service"),
    ("ms_workflow",      "execution/workflow-service"),
    ("ms_verification",  "verification/verification-service"),
    ("ms_evidence",      "verification/evidence-service"),
    ("ms_gps",           "verification/gps-service"),
    ("ms_audit",         "verification/audit-service"),
    ("ms_media",         "verification/media-service"),
    ("ms_submission",    "verification/submission-service"),
    ("ms_graph",         "graph/graph-service"),
    ("ms_reputation",    "reputation/portfolio-service"),
    ("ms_impact_score",  "reputation/impact-score-service"),
    ("ms_org_score",     "reputation/organization-score-service"),
    ("ms_marketplace",   "marketplace/opportunity-service"),
    ("ms_matching",      "marketplace/matching-service"),
    ("ms_report",        "reporting/report-service"),
    ("ms_analytics",     "reporting/analytics-service"),
    ("ms_forecast",      "reporting/forecast-service"),
    ("ms_notification",  "notifications/notification-service"),
    ("ms_chat",          "chat-service"),
    ("ms_billing",       "billing/billing-service"),
    ("ms_payment",       "billing/payment-service"),
    ("ms_subscription",  "billing/subscription-service"),
    ("ms_partner",       "partners/partner-service"),
    ("ms_coin",          "coin/coin-service"),
    ("ms_ads",           "ads/ads-service"),
    ("ms_platform",      "platform/platform-service"),
]
\end{lstlisting}

\subsection{Package Registration Function}

Each service is dynamically registered as a top-level Python package. This allows standard Python imports to resolve service modules correctly:

\begin{lstlisting}[language=Python, caption=Dynamic Package Registration]
def _register_service_package(alias: str, rel: str) -> None:
    \"\"\"Register backend/<rel> as a top-level package named alias.\"\"\"
    svc_root = os.path.join(BACKEND, rel)
    pkg = types.ModuleType(alias)
    pkg.__path__ = [svc_root]
    pkg.__package__ = alias
    sys.modules[alias] = pkg
\end{lstlisting}

\subsection{Duplicate Table Resolution}

A critical challenge of merging 30+ independent services is that many define the same SQLAlchemy ORM tables (e.g., \texttt{users}, \texttt{organizations}). The solution is a monkey-patch on \texttt{Table.\_\_new\_\_} that deduplicates table registrations:

\begin{lstlisting}[language=Python, caption=SQLAlchemy Table Deduplication]
import sqlalchemy as _sa
from sqlalchemy import Table as _Table

_orig_table_new = _Table.__new__

def _patched_table_new(cls, *args, **kwargs):
    metadata = kwargs.get("metadata") or \
               (args[1] if len(args) >= 2 else None)
    name = kwargs.get("name") or \
           (args[0] if args else None)
    schema = kwargs.get("schema")

    if metadata is not None and name is not None:
        key = (schema, name) if schema else name
        if key in metadata.tables:
            return metadata.tables[key]

    return _orig_table_new(cls, *args, **kwargs)

_Table.__new__ = staticmethod(_patched_table_new)
\end{lstlisting}

\subsection{Route Merging}

After registering packages and patching tables, the combined backend iterates through all services and merges their FastAPI routes into a single application:

\begin{lstlisting}[language=Python, caption=Route Merging]
from fastapi import FastAPI
from fastapi.middleware.cors import CORSMiddleware

app = FastAPI(
    title="MASAREN Combined Backend",
    version="0.1.0"
)

for alias, rel in SERVICES:
    try:
        _register_service_package(alias, rel)
        mod = importlib.import_module(f"{alias}.src.main")
        svc = getattr(mod, "app", None)

        if svc is None:
            continue

        for r in list(svc.routes):
            if getattr(r, "path", None) in SKIP:
                svc.routes.remove(r)

        app.include_router(svc.router)

        for exc_type, handler in \
                svc.exception_handlers.items():
            if exc_type not in app.exception_handlers:
                app.add_exception_handler(exc_type, handler)

        _included.append(alias)
        print(f"[combined] OK included {alias}", flush=True)
    except Exception as e:
        print(f"[combined] FAIL {alias}: {e!r}", flush=True)

app.add_middleware(
    CORSMiddleware,
    allow_origins=os.getenv(
        "CORS_ORIGINS",
        "http://localhost:3000,http://localhost:3001"
    ).split(","),
    allow_credentials=True,
    allow_methods=["*"],
    allow_headers=["*"],
)
\end{lstlisting}

\subsection{Service Skip List}

Certain routes are excluded from the combined backend to avoid conflicts:
\begin{itemize}
  \item Gateway routes (\texttt{/api/v1/*}) -- the gateway should not proxy to itself
  \item Duplicate health check endpoints (only one \texttt{/health} is kept)
  \item OpenAPI documentation endpoints (only the combined app's docs are exposed)
\end{itemize}

\clearpage

% ============================================================
% SECTION 6: WEB APPLICATION ARCHITECTURE
% ============================================================
\section{Web Application Architecture}

\subsection{Component Architecture}

The web application at \texttt{apps/web} is a React 18.2 + TypeScript application built with Vite 5.1. It follows a component-based architecture with clear separation of concerns:

\subsubsection{Directory Structure}

\begin{verbatim}
apps/web/src/
  components/
    layout/          TopNav, Sidebar, Header, RightRail
    common/          Branding, ErrorBoundary, Icons
  pages/             30+ page components
  services/          API client layer
  store/             React Context providers
  hooks/             Custom React hooks
  types/             TypeScript type definitions
  test/              Test utilities and setup
  i18n.tsx           Internationalization configuration
\end{verbatim}

\subsubsection{Page Structure}

The application defines 30+ pages accessible through React Router v6:

\begin{longtable}{|p{3.5cm}|p{3cm}|p{6cm}|}
\hline
\textbf{Page} & \textbf{Route} & \textbf{Description} \\
\hline
\endhead
LandingPage & /, /landing & Public landing with platform overview \\
\hline
LoginPage & /login & Authentication with email/password and Google OAuth \\
\hline
RegisterPage & /register & User registration with email verification \\
\hline
DashboardPage & /dashboard & Personal dashboard with metrics and activity \\
\hline
ProgramsPage & /programs & Program discovery and enrollment \\
\hline
TasksPage & /tasks & Task management and assignment \\
\hline
EvidencePage & /evidence & Evidence submission and tracking \\
\hline
PassportPage & /passport & Digital impact passport \\
\hline
OpportunitiesPage & /opportunities & Marketplace opportunity discovery \\
\hline
PortfolioPage & /portfolio & Impact asset portfolio \\
\hline
ProfilePage & /profile & User profile management \\
\hline
PersonnelProfilePage & /personnel/:id & Public personnel profile \\
\hline
InstitutionProfilePage & /institution/:id & Institution profile view \\
\hline
AdminPage & /admin & Administrative dashboard (ADMIN/SUPER\_ADMIN) \\
\hline
NotificationsPage & /notifications & Notification center \\
\hline
ReportsPage & /reports & Analytics and reports \\
\hline
BillingPage & /billing & Billing and subscription management \\
\hline
ChatPage & /chat & Real-time messaging \\
\hline
FeedPage & /feed & Social activity feed \\
\hline
WalletPage & /wallet & MRC wallet management \\
\hline
\ldots & \ldots & Additional pages for ads, payments, revenue, challenges, etc. \\
\hline
\end{longtable}

\subsection{Routing Architecture}

The application implements protected routes with role-based access control:

\begin{lstlisting}[language=TypeScript, caption=Protected Route Implementation]
function RequireRole({
  roles,
  children
}: {
  roles: string[];
  children: ReactElement;
}) {
  const { user, isAuthenticated } = useAuth();
  if (!isAuthenticated) return <Navigate to="/login" replace />;
  if (user && !roles.includes(user.role))
    return <Navigate to="/dashboard" replace />;
  return children;
}
\end{lstlisting}

\subsection{State Management}

The application uses a dual state management approach:

\begin{description}
  \item[React Context] Client-side state for authentication (\texttt{AuthProvider}), theme preferences, and UI state
  \item[TanStack Query] Server state management with automatic caching, background refetching, and optimistic updates
\end{description}

\subsection{Internationalization}

The application supports bilingual English/Arabic operation:

\begin{lstlisting}[language=TypeScript, caption=i18n Configuration]
// apps/web/src/i18n.tsx
import i18n from 'i18next';
import { initReactI18next } from 'react-i18next';

i18n.use(initReactI18next).init({
  resources: {
    en: { translation: enTranslations },
    ar: { translation: arTranslations }
  },
  lng: 'en',
  fallbackLng: 'en',
  interpolation: { escapeValue: false }
});
\end{lstlisting}

\subsection{Authentication Flow}

\begin{enumerate}
  \item User navigates to login page
  \item User authenticates via email/password or Google OAuth popup
  \item Backend issues JWT access token and refresh token
  \item Axios interceptor attaches JWT to all API requests
  \item On 401 response, interceptor attempts token refresh
  \item On refresh failure, user is redirected to login
\end{enumerate}

\subsection{API Client Layer}

The services layer provides type-safe API interaction:

\begin{lstlisting}[language=TypeScript, caption=API Client Structure]
// apps/web/src/services/
import axios from 'axios';

const api = axios.create({
  baseURL: import.meta.env.VITE_API_URL ||
           'http://localhost:8080/api/v1',
});

api.interceptors.request.use((config) => {
  const token = getAccessToken();
  if (token) {
    config.headers.Authorization = `Bearer ${token}`;
  }
  return config;
});

api.interceptors.response.use(
  (response) => response,
  async (error) => {
    if (error.response?.status === 401) {
      await refreshToken();
      return api.request(error.config);
    }
    return Promise.reject(error);
  }
);
\end{lstlisting}

\clearpage

% ============================================================
% SECTION 7: MOBILE APPLICATION ARCHITECTURE
% ============================================================
\section{Mobile Application Architecture}

\subsection{Architecture Overview}

The mobile application at \texttt{apps/mobile} is built with Flutter 3.22+ and follows the MVVM (Model-View-ViewModel) architecture pattern with Provider for state management. The app provides the same feature set as the web application, optimized for mobile form factors with offline support and device-native capabilities.

\subsection{Directory Structure}

\begin{verbatim}
apps/mobile/lib/
  features/            Feature-specific modules
  providers/           ChangeNotifier providers
  services/            API client and service layer
  widgets/             Reusable UI components
  models/              Data models and DTOs
  screens/             Full-page screen widgets
  utils/               Utility functions and helpers
  main.dart            Application entry point
\end{verbatim}

\subsection{Screen Structure}

The Flutter app implements the following screens:
\begin{itemize}
  \item Login and Registration screens with Google OAuth integration
  \item Dashboard with activity summary and quick actions
  \item Program discovery and enrollment screens
  \item Task list and detail screens with mobile-optimized forms
  \item Evidence capture screen with camera integration and GPS binding
  \item Digital passport visualization
  \item Opportunity marketplace with search and filters
  \item Impact portfolio with charts and metrics
  \item Profile management with edit capabilities
  \item Real-time chat interface
  \item Social feed with infinite scroll
  \item Notification center with push notification handling
  \item Billing and subscription management
  \item Wallet for MRC balance and transactions
\end{itemize}

\subsection{State Management}

The app uses the Provider package for state management:

\begin{lstlisting}[language=Dart, caption=Provider Setup]
// main.dart
void main() {
  runApp(
    MultiProvider(
      providers: [
        ChangeNotifierProvider(create: (_) => AuthProvider()),
        ChangeNotifierProvider(create: (_) => ThemeProvider()),
        ChangeNotifierProvider(create: (_) => ProgramProvider()),
        ChangeNotifierProvider(create: (_) => TaskProvider()),
        ChangeNotifierProvider(create: (_) => EvidenceProvider()),
        ChangeNotifierProvider(create: (_) => NotificationProvider()),
      ],
      child: const MasarenApp(),
    ),
  );
}
\end{lstlisting}

\subsection{Navigation Architecture}

The app uses Flutter Navigator with a custom bottom navigation bar:

\begin{itemize}
  \item \textbf{Bottom Navigation Tabs:} Dashboard, Programs, Tasks, Profile, More
  \item \textbf{Drawer Menu:} Settings, Help, About, Logout
  \item \textbf{Deep Linking:} Notification-driven navigation to specific screens
  \item \textbf{Modal Routes:} Evidence capture, chat, and form modals
\end{itemize}

\subsection{Mobile-Specific Features}

The mobile app leverages device capabilities not available on the web:

\begin{itemize}
  \item \textbf{Camera Integration:} Direct evidence capture through device camera
  \item \textbf{GPS Location Binding:} Automatic coordinates attachment to submissions
  \item \textbf{Offline Mode:} Queue submissions when offline, sync when connected
  \item \textbf{Push Notifications:} FCM-based real-time alerts for assignments and verification
  \item \textbf{Biometric Auth:} Fingerprint/Face ID for quick app unlock
  \item \textbf{Image Gallery:} Photo library access for evidence selection
  \item \textbf{File Management:} Download and open impact certificates and reports
\end{itemize}

\subsection{API Integration}

\begin{lstlisting}[language=Dart, caption=API Service Layer]
class ApiService {
  static const String baseUrl =
    'http://localhost:8080/api/v1';

  Future<Map<String, dynamic>> get(
    String endpoint, {
    Map<String, String>? headers,
  }) async {
    final response = await http.get(
      Uri.parse('$baseUrl/$endpoint'),
      headers: {
        'Content-Type': 'application/json',
        'Authorization': 'Bearer ${await getToken()}',
        ...?headers,
      },
    );
    return json.decode(response.body);
  }

  Future<Map<String, dynamic>> post(
    String endpoint, {
    Map<String, dynamic>? body,
  }) async {
    final response = await http.post(
      Uri.parse('$baseUrl/$endpoint'),
      headers: {
        'Content-Type': 'application/json',
        'Authorization': 'Bearer ${await getToken()}',
      },
      body: json.encode(body),
    );
    return json.decode(response.body);
  }
}
\end{lstlisting}

\subsection{Bilingual Support}

The mobile app provides full bilingual support similar to the web:

\begin{lstlisting}[language=Dart, caption=Internationalization Setup]
class MasarenLocalizations {
  static const Map<String, Map<String, String>> _localized = {
    'en': {
      'app.title': 'MASAREN',
      'nav.home': 'Home',
      'nav.programs': 'Programs',
      'nav.tasks': 'Tasks',
      'nav.profile': 'Profile',
    },
    'ar': {
      'app.title': 'مسارين',
      'nav.home': 'الرئيسية',
      'nav.programs': 'البرامج',
      'nav.tasks': 'المهام',
      'nav.profile': 'الملف الشخصي',
    },
  };
}
\end{lstlisting}

\clearpage
